\section{Endpoint distributions and conditional recovery}

The probabilities assigned to two atom laws determine which combination
of their endpoints can be recovered. We first weight the laws by the
actual endpoint norms to control their average. We then compare signed
and conditional recovery when finite laws only approximate the endpoints,
and finally keep the original radius and a small center tail in the
budget.

\subsection{Weighting the endpoint distributions by their norms}
\label{molw:sec}

The two actual endpoint norms determine a probability law that controls
their average directly. We work in the infinite-height coordinate
predual $B_\infty$ of Section~\ref{binf:definition}. Its atomic gauge
and finite-coordinate polar descriptions agree by
Lemma~\ref{binf:finite-hull}; the completion is the norm-closed coordinate span
specified there, as distinguished in Remark~\ref{presentation:infinite-height}.
We use its represented atoms $m=\sum_i a_i d_{S_i}$, its ancestral
projections $P_A$ and $Q_A=I-P_A$, and the cone estimates from
Lemmas~\ref{binf:projections} and~\ref{binf:cones}.

\begin{theorem}\label{molw:average}
Let $A\subset T_\infty$ be finite and ancestral, let
$x\in B_\infty$ be supported in $A$, and let $z\in\ker P_A$. Write
\[
 r=\|x\|,\qquad L_+=\|x+z\|,\qquad L_-=\|x-z\|,
 \qquad L=\frac{L_++L_-}{2}.
\]
Then $L_+,L_-\ge r$ and
\begin{equation}\label{molw:lens}
 \|z\|\le C\sqrt{L(L-r)},\qquad C=\sqrt2+\sqrt6.
\end{equation}
\end{theorem}

\begin{proof}
Projection onto $A$ gives $L_+,L_-\ge r$. Thus $L<r$ is impossible.
If $r=0$, then $x=0$ and $L_+=L_-=L=\|z\|$, and
\eqref{molw:lens} follows because $C\ge1$. This includes $L=0$.
If an endpoint norm is zero, projection forces $r=0$ and then $z=0$,
so no zero denominator is needed in what follows. For $r>0$, divide
$x$ and $z$ by $r$; it is enough to prove the assertion with $r=1$.
In this case $L_+,L_-\ge1$.

First suppose $z$ is finitely supported, and choose a finite ancestral
set $B$ containing $A$ and the support of $z$. By
Lemma~\ref{binf:projections}, choose a head test $f=P_Af$ with
$\rho(f)=1$ and $x(f)=1$. The segment test norm gives
\begin{equation}\label{molw:test}
 \sum_i f(S_i)^2\le1
\end{equation}
for every disjoint segment family.

Choose a sign $\sigma$ with probabilities
\begin{equation}\label{molw:weights}
 \mathbb P\{\sigma=1\}=\frac{L_+}{2L},\qquad
 \mathbb P\{\sigma=-1\}=\frac{L_-}{2L}.
\end{equation}
Conditionally on that sign, choose a finite distribution of represented
atoms with mean $(x+\sigma z)/L_\sigma$. Such distributions exist by
Lemma~\ref{binf:finite-hull}; the random data retain a chosen
coefficient list and segment family for every atom. These probabilities
give both identities
\begin{equation}\label{molw:means}
 \mathbb E m=\frac{x}{L},\qquad
 \mathbb E(\sigma m)=\frac{z}{L}.
\end{equation}
The first identity cancels the tail. The second recovers the desired
displacement without passing to one conditional distribution.

Apply Corollary~\ref{stab:bounded-multiplier} to this finite law,
with multiplier $\nu=\sigma$. Its mean $x/L$ is supported in $A$,
so the center-tail parameter in that corollary is zero. Its evaluation
deficit under $f$ is
\[
 \delta=1-(x/L)(f)=1-1/L\ge0.
\]
The first inequality in \eqref{stab:bounded-multiplier-bound}, together
with the two identities in \eqref{molw:means}, therefore gives
\[
 \frac{\|z\|}{L}
 =\left\|Q_A\mathbb E(\sigma m)\right\|
 \le\sqrt{6\delta}+\sqrt{2\delta}
 =(\sqrt6+\sqrt2)\sqrt{1-1/L}.
\]
Multiplying by $L$ proves \eqref{molw:lens} when $r=1$.
The multiplier may depend on the represented atom; the corollary requires
only $|\sigma|\le1$, which holds for the endpoint sign used here.

For general $z\in\ker P_A$, use the finite tail approximants from
Lemma~\ref{binf:projections}. Their endpoint norms and averages
converge, so the finite-support inequality passes to the limit.
Finally, rescaling from $r=1$ gives
$rC\sqrt{(L/r)(L/r-1)}=C\sqrt{L(L-r)}$.
\end{proof}

\begin{corollary}\label{molw:modulus}
For every $t>0$, put
\begin{equation}\label{molw:rho}
 \omega(t)=\frac{\sqrt{1+4(t/C)^2}-1}{2},\qquad C=\sqrt2+\sqrt6.
\end{equation}
Then the averaged asymptotic modulus of $B_\infty$ satisfies
$\widehat\delta_{B_\infty}(t)\ge\omega(t)/2>0$. The bound holds when
the cofinite infimum allows every perturbation direction of norm at
least one.
\end{corollary}

\begin{proof}
For a finite-support unit center $x$, choose a finite ancestral set $A$
containing its support and take $F=\ker P_A$. If $y\in F$ and
$\|y\|\ge1$, Theorem~\ref{molw:average} applied to $z=ty$ yields
$t^2\le C^2 L(L-1)$. Since $L\ge1$, solving this quadratic gives
$L-1\ge\omega(t)$. The projection $P_A$ has finite rank, so $F$ is
closed and has finite codimension.

For an arbitrary unit center $x$, density and normalization provide
a finite-support unit vector $x_0$ with
$\|x-x_0\|<\omega(t)/2$. Choose the subspace associated with $x_0$
before choosing $y$. Each endpoint norm changes by at most
$\|x-x_0\|$, uniformly for all such $y$, so their average exceeds
$1+\omega(t)/2$. Taking the infimum over $y$, the supremum over
$F$, and then the infimum over unit centers gives the assertion.
This argument imposes no upper bound on $t$.
\end{proof}

We now return to the finite-height forest space $X=J^*$ of
Proposition~\ref{presentation:identity}.

\subsection{Signed and conditional recovery from finite approximations}
\label{mola:conditional-section}\label{mol:closed-hull-signed-section}

Finite atom laws approximate completed endpoints in norm. Their mixture
therefore has a small tail even when the intended center is supported
on the head. The next two estimates keep this approximation error
explicit. Recovering one conditional law incurs a factor of two in its
absolute boundary mass; recovering the signed mixture avoids that factor.
The conditional estimate allows the displacement to have a nonzero head,
whereas the signed statement below is formulated for a pure tail.

\begin{proposition}\label{mola:conditional}
Let $H$ be finite and ancestor-closed, $x=P_Hx$, and
$x+h,x-h\in B_X$, with arbitrary $h\in X$.
Suppose $0<d<1$ and $\|x\|\ge1-d$. For every $\eta>0$, choose
finite represented-atom laws with means $z_\pm$ such that
$\|z_\pm-(x\pm h)\|\le\eta$. Such laws exist, and every such pair gives
\begin{equation}\label{mola:conditional-prelimit}
 \|Q_Hh\|\le\eta+2\sqrt{d+\eta}
       +2\sqrt{(\eta+\sqrt{2(d+\eta)})^2+4(d+\eta)}.
\end{equation}
Consequently,
\[
 \|Q_Hh\|\le(2+2\sqrt6)\sqrt d.
\]
\end{proposition}
\begin{proof}
By Lemma~\ref{lem:finite-hull}, choose $p=P_Hp$ with
$\rho(p)=1$ and $x(p)=\|x\|\ge1-d$.
The two ball inclusions give $x(p)\pm h(p)\le1$, and hence
$|h(p)|\le d$.  Fix $\eta>0$.  Corollary~\ref{cor:closed-atomic-ball}
provides finite represented-atom mixtures approximating the two
endpoints within $\eta$.  Write
$\mathbb E_+$, $\mathbb E_-$ for their expectations and
$\mathbb E=(\mathbb E_++\mathbb E_-)/2$ for the equal mixture.
For each atom $Z=\sum_i a_i d_{S_i}$ set
$b_i=p(S_i)$, $\widehat Z=\sum_i b_i d_{S_i}$,
$D=\sum_i(a_i-b_i)^2$, and $L=1-\sum_i b_i^2$.
Pointwise $D+L\le2(1-Z(p))$.
The mixture pairs with $p$ by at least $1-d-\eta$, whereas each
conditional law pairs with $p$ by at least $1-2d-\eta$.  Therefore
\[
 \mathbb ED,\mathbb EL\le2(d+\eta),\qquad
 \mathbb E_\pm D,\mathbb E_\pm L\le4d+2\eta\le4(d+\eta).
\]
In particular the conditional replacement error is at most
$2\sqrt{d+\eta}$.

Write $z_\pm=\mathbb E_\pm Z$ and $z=(z_++z_-)/2$.
These finite atom means lie in $B_X$. Their approximation to the
intended endpoints gives
\[
 \|Q_Hz\|\le\eta,\qquad
 1-z(p)\le d+\eta,\qquad 1-z_\pm(p)\le2d+\eta.
\]
Apply Theorem~\ref{stab:weighted} to $z_+,z_-$ with equal weights
and this same head test. It gives
\[
 \|Q_Hz_\pm\|
 \le2\sqrt{(\eta+\sqrt{2(d+\eta)})^2+4(d+\eta)}
       +\sqrt{4d+2\eta}.
\]
The conditional boundary-mass factor is two, while the conditional
replacement term is at most $2\sqrt{d+\eta}$. Restoring the endpoint
approximation costs at most $\eta$, since $Q_Hx=0$ and $Q_H$ is
contractive. This proves \eqref{mola:conditional-prelimit}.
Letting $\eta\downarrow0$ gives the asserted limit.
\end{proof}
For signed recovery, retain the endpoint sign in the equal mixture of
the two approximating laws. The following estimate shows exactly which
conditional loss disappears. The head again remains fixed, and the
displacement is not projected to finite support.

\begin{proposition}\label{mol:closed-hull-signed}
In the forest space $X=J^*$, let $H$ be finite and ancestral,
$x=P_Hx$, $y=Q_Hy$, and $x\pm y\in B_X$. Put $d=1-\|x\|$.
For every $\varepsilon>0$, the finite atomic approximation argument gives
\begin{equation}\label{mol:closed-hull-error}
 \|y\|\le e_\varepsilon+
 \sqrt{e_\varepsilon^2+4(d+\varepsilon)},\qquad
 e_\varepsilon=\varepsilon+\sqrt{2(d+\varepsilon)}.
\end{equation}
Consequently $\|y\|\le(\sqrt2+\sqrt6)\sqrt d$.
\end{proposition}

\begin{proof}
Choose a head test $p=P_Hp$ with $\rho(p)\le1$ and
$x(p)=\|x\|$. Approximate $x\pm y$ within $\varepsilon$ by
finite represented-atom means $z_\pm$, using
Corollary~\ref{cor:closed-atomic-ball}. On their equal mixture, let
$\sigma=1$ on the plus law and $\sigma=-1$ on the minus law.
Then $\|\mathbb EA-x\|\le\varepsilon$ and
$\|\mathbb E(\sigma A)-y\|\le\varepsilon$.
Its mean test deficit is at most $d+\varepsilon$, and its mean
tail has norm at most $\varepsilon$. Corollary~\ref{stab:bounded-multiplier}
applies to this finite law with multiplier $\sigma$, and hence gives
\[
 \left\|Q_H\mathbb E(\sigma A)\right\|
 \le\sqrt{(\varepsilon+\sqrt{2(d+\varepsilon)})^2
                   +4(d+\varepsilon)}
       +\sqrt{2(d+\varepsilon)}.
\]
There is no conditional boundary-mass factor: the signed expectation
recovers the difference of the two approximating endpoints directly.
Since $Q_Hy=y$, adding the endpoint approximation error
$\varepsilon$ proves \eqref{mol:closed-hull-error} with the stated
$e_\varepsilon$.
Let $\varepsilon$ decrease to zero. For endpoints of norm at most
$1+\lambda$ about a supported unit center, first divide by
$1+\lambda$ and take $d=\lambda/(1+\lambda)$ in the same calculation.
\end{proof}

For a general near-unit center, first approximate it by its projection
onto one finite ancestor-closed head, and divide that projected center
and each displacement by one plus the center-approximation error.
Both endpoint bounds then hold in the unit ball, so
Proposition~\ref{mola:conditional} controls every displacement tail
using that same head. For an infinite family $(v_j)$ with
$x\pm v_j\in B_X$, the bounded head projections $(P_Hv_j)$ are
totally bounded. Thus, if the family is $\varepsilon$-separated,
the center satisfies $\|x\|\le1-\delta$ for some $\delta>0$
depending only on $\varepsilon$. This is the conclusion of
Corollary~\ref{mola:separated-sequential}.

\subsection{Endpoint laws at their original radius}
\label{molb:unscaled-section}

We work in the finite-height forest space $X$ of
Proposition~\ref{presentation:identity}, with primal segment norm
$\rho$ on $J$. Keeping the endpoint radius in the represented
coefficients gives a joint budget for their replacement error and
the unused squared head test.

\begin{proposition}\label{molb:unscaled-tail}
Let $\delta>0$ and $x\in X$ satisfy $\|x\|=1$. There are a finite
ancestral head $H$ and a test $p=P_Hp\in J$ such that
\begin{equation}\label{molb:unscaled-head}
 \rho(p)\le1,\qquad x(p)\ge1-\delta,\qquad
 \|Q_Hx\|\le\delta.
\end{equation}
For this same head and every $y\in X$ satisfying
$\|x+y\|,\|x-y\|\le1+\delta$, put $h=6\delta+\delta^2$. Then
\begin{equation}\label{molb:unscaled-bound}
 \max\{\|Q_H(x+y)\|,\|Q_H(x-y)\|,\|Q_Hy\|\}
 \le F(\delta):=
 \delta+\sqrt{2h}+2\sqrt{(2\delta+\sqrt h)^2+2h}.
\end{equation}
In particular, for every $\varepsilon>0$ one can choose $\delta>0$
independently of $x$ so that this head satisfies
$\|Q_Hy\|<\varepsilon$ for every such $y$.
\end{proposition}

\begin{proof}
Coordinate density and contractivity give a finite ancestral $H$
with $\|Q_Hx\|\le\delta$. Under the identification $X^*=J$ in
Proposition~\ref{presentation:identity}, choose a norming test
$f\in J$ with $\rho(f)=1$ and $x(f)=1$. Set $p=P_Hf$.
Then $\rho(p)\le1$ and
$x(p)=1-(Q_Hx)(f)\ge1-\delta$, proving
\eqref{molb:unscaled-head}. This choice uses only $x$ and $\delta$.

Fix any permitted $y$. The radius-$(1+\delta)$ version of
Corollary~\ref{cor:closed-atomic-ball} provides two finite laws of
represented functionals
\[
 V=\sum_i a_i d_{S_i},\qquad
 S_i\text{ pairwise disjoint},\qquad
 \sum_i a_i^2\le(1+\delta)^2,
\]
whose respective means approximate $x+y$ and $x-y$ within $\delta$.
Write $\mathbb E_+$ and $\mathbb E_-$ for these laws and
$\mathbb E=(\mathbb E_++\mathbb E_-)/2$ for their equal mixture.
Thus
\[
 \|\mathbb E_\pm V-(x\pm y)\|\le\delta,
 \qquad \|\mathbb EV-x\|\le\delta,
 \qquad (\mathbb EV)(p)\ge1-2\delta.
\]
The last bound uses $\rho(p)\le1$ and
\eqref{molb:unscaled-head}.

In each represented outcome define
\[
 b_i=p(S_i),\qquad W=\sum_i b_i d_{S_i},\qquad
 L=1-\sum_i b_i^2,\qquad E=\sum_i(a_i-b_i)^2.
\]
Both $L$ and $E$ are nonnegative. Expanding the square and taking
expectations gives the unscaled joint budget
\begin{equation}\label{molb:unscaled-budget}
 \mathbb E(E+L)
 =\mathbb E\left(1+\sum_i a_i^2-2V(p)\right)
 \le(1+\delta)^2+1-2(1-2\delta)
 =6\delta+\delta^2=h.
\end{equation}
The disjoint-segment coefficient estimate
\eqref{stab:sign:coefficients} gives $\|V-W\|\le\sqrt E$.
Cauchy--Schwarz, nonnegativity, and the equal mixture therefore yield
\begin{equation}\label{molb:unscaled-replacement}
 \mathbb E\|V-W\|\le\sqrt h,
 \qquad
 \mathbb E_\pm\|V-W\|
 \le\sqrt{\mathbb E_\pm E}\le\sqrt{2h}.
\end{equation}

Let $G_H$ be the first exterior vertices. A segment wholly outside
$H$ has replacement coefficient zero. At each $g\in G_H$ there is
at most one crossing segment in an outcome; let $q_g$ be its head
sum $b_i$, or zero if it is absent. Thus $W(e_g)=q_g$ and the
restriction of $Q_HW$ to the cone at $g$ is either zero or $q_g$
times a segment functional starting at $g$. Put
$m_g=\mathbb E|q_g|$ and $v_g=\mathbb E q_g$.
These vectors have finite support. Antichain contraction, followed
by \eqref{molb:unscaled-replacement}, gives
\begin{equation}\label{molb:unscaled-signed}
 \|v\|_2\le\|Q_H\mathbb EW\|
 \le\|Q_Hx\|+\|\mathbb EV-x\|
       +\|\mathbb E(W-V)\|
 \le2\delta+\sqrt h.
\end{equation}

Apply Lemma~\ref{lem:opposite-crossing} to two independent copies
of the represented segment families, both tested against $p$.
Its simultaneous two-piece splits retain the longer segment's
exterior endpoint, and their available squared-test budgets are
$L$ and $L'$. Consequently
\[
 2\sum_{g:q_gq'_g<0}|q_gq'_g|\le L+L'.
\]
Independence factors every product expectation. In combination with
\eqref{molb:unscaled-budget}, this gives
\begin{equation}\label{molb:unscaled-unsigned}
 \|m\|_2^2-\|v\|_2^2
 =2\mathbb E\sum_{g:q_gq'_g<0}|q_gq'_g|
 \le2\mathbb EL\le2h,
 \qquad
 \|m\|_2\le\sqrt{(2\delta+\sqrt h)^2+2h}.
\end{equation}

For either conditional law, the absolute mass in the cone at $g$
is at most $\mathbb E_\pm|q_g|\le2m_g$. The rooted-mixture
estimate in Lemma~\ref{stab:coef:tail-assembly}, equivalently the
subtree decomposition of Proposition~\ref{presentation:subtrees},
therefore gives
\begin{equation}\label{molb:unscaled-conditional}
 \|Q_H\mathbb E_\pm W\|\le2\|m\|_2.
\end{equation}
Restoring the conditional coefficient error and the endpoint
approximation yields
\[
 \|Q_H(x\pm y)\|
 \le\delta+\sqrt{2h}+2\|m\|_2
 \le F(\delta).
\]
Taking half the difference of the two endpoint tails proves the
same bound for $Q_Hy$. Finally $F(\delta)\to0$ as
$\delta\downarrow0$. Choosing $F(\delta)<\varepsilon$ before
choosing $x$, and then choosing $H$ as above before $y$, proves
the uniform assertion.
\end{proof}
